\section{Introduction}
\label{sec:introduction}

Multiple-changepoint analysis represents an ordered series by contiguous segments
within which a sampling-model parameter is constant. It is used in statistical
signal processing, genomics, environmental science, and econometrics. The
statistical task is to estimate the segment count, changepoint locations,
segment parameters, and their uncertainty. Bayesian methods assign posterior
probabilities to segmentations and can average over these alternatives when
estimating a piecewise-constant signal or predicting future observations
\citep{barry1992ppm,barry1993bayesCP,fearnhead2006exact,hutter2006bpcr}.

For an offline changepoint model with independent segment parameters and a
factorizing partition prior, the probability of a partition is a product of
integrated segment likelihoods and local prior terms. Dynamic programming
then sums these products without enumerating every partition
\citep{barry1992ppm,fearnhead2006exact,hutter2006bpcr}. This formulation is
well established. Its practical use across heterogeneous outcomes depends on
calculating segment marginal likelihoods with the correct sampling-model
normalization and integrating the partition prior with the same care.

\BayesBreak develops a generalized hierarchical formulation of this model.
For regular exponential-family observations with proper conjugate priors,
segment marginal likelihoods and posterior moments follow from cumulative
sufficient statistics. The resulting partition recursions do not depend on
which member of the family supplied the segment scores. The framework includes
nonuniform partitions on irregular physical coordinates, several sequences
sharing changepoints, known groups, and a finite latent-group template
criterion. Nonconjugate segment integrals can also be approximated locally,
but exact posterior interpretation of the resulting dynamic program requires
control of the segment-level approximation error.

\paragraph{Contributions.}
The contributions of this paper are as follows.
\begin{enumerate}
\item We specify a common conjugate segment calculation for exponential-family
  sampling models with known exposure, trial-count, and precision descriptors,
  explicitly distinguishing ordinary marginal likelihoods from powered
  likelihood integrals. The calculation provides segment marginal likelihoods
  and the posterior moments required for Bayesian piecewise-constant regression.
\item We derive normalized sum-product recursions for the segment-count
  posterior, ordered and unordered changepoint marginals, and posterior signal
  moments. We also derive a separate max-sum recursion for joint MAP
  segmentation, with immediate proofs of the identities and decision rules.
\item We extend the local-factor model to physical segment lengths and
  candidate-boundary weights, and to the exact pooling of independently sampled
  sequences under a common segmentation.
\item We study known-group pooling and a finite latent-template criterion for
  unknown groups. A Jensen minorization leads to alternating allocation and
  responsibility-weighted optimal-partition updates, with monotonicity for the
  specified objective rather than a claim of full Bayesian mixture inference.
\item We analyze the effect of segment marginal-likelihood approximation error
  on partition probabilities and posterior odds. The analysis is conditional
  on a uniform bound across admissible segments, not on a numerical method's
  name or an unverified finite-sample error estimate.
\end{enumerate}

The experiments examine simulation recovery, posterior calibration, numerical
approximations, computation time, and applications to well-log measurements,
array-CGH, and methylation. Existing results are reported
separately from the additional experiments specified in
Section~\ref{sec:experimental-design}. The latter are planned evaluations and
are not used to support conclusions about accuracy or comparative performance.

Section~\ref{sec:related} reviews the relevant statistical literature.
Section~\ref{sec:setup} defines the sampling model, partition, and inferential
targets. Sections~\ref{sec:model} and~\ref{sec:exact-inference} derive segment
integration and exact partition inference. Section~\ref{sec:extensions}
covers hierarchical extensions. Sections~\ref{sec:approx-predict} and
\ref{sec:algorithms} address numerical approximation, prediction, and
implementation. The experimental protocols and results follow in
Sections~\ref{sec:experimental-design} and~\ref{sec:results}.
